\section{章节罗盘：帝国终结前的政治基础}
\label{sec:revolution-compass}

1911年的革命不是突然脱离晚清改革的一页。本章从咨议局、商会、报刊、新军和铁路公司形成的政治场所写起，考察它们怎样扩大组织与发言的机会，也怎样受财产、教育、地域和军费的限制。国有化与外债争议之所以能引发保路运动，正在于地方权利、中央债务和帝国的财政集中被压在同一条线路上。

武昌之后的起事、各省独立、议和和退位将被视为互相牵动而不同步的过程。问题不只是谁取代皇帝，更是谁能承继税源、军队、债务和地方日常秩序。本章的结尾会将革命放回全卷的长线：征服帝国留下的边疆、财政和军事协调机制，在共和国名义出现后仍须重新谈判。
